\documentclass{ximera}
\title{Tensors on manifolds}

\begin{document}

    \begin{abstract}  Tensors on manifolds.   \end{abstract}
    \maketitle

 A reference for this material is Chapter 12  of John M. Lee. Introduction to smooth manifolds. Second edition. Grad. Texts in Math., 218. Springer, New York, 2013. xvi+708pp. ISBN: 978-1-4419-9981-8. 


\begin{definition}
    The $(k,\ell)$\emph{-tensor bundle} over $M$ is defined as
    \[    T ^{k}_{\ell} (M) : = \underbrace{ T^{\ast}M \otimes \cdots \otimes T^{\ast}M }_{k  \text{ times}} \otimes  \underbrace{ TM \otimes \cdots \otimes TM }_{\ell \text{ times}} .   
    \]
    We also denote
    \[   \mathcal{T}^k_{\ell} (M) : = \Gamma ( T ^k_{\ell} (M) ).                 \]
    A section in $\mathcal{T}^k_{\ell} (M)$ is called a \emph{tensor field of type} $(\ell ,  k )$. 
\end{definition}
Recall that the modules of sections are finitely-generated projective, so we have a natural isomorphism
\[   \mathcal{T}^k_{\ell} (M) \cong Hom_{C^{\infty}(M)}  (  \mathfrak{X} (M ) ^{\otimes k} , \mathfrak{X} (M) ^{ \otimes  \ell}   )   .                 \]
In other words, a tensor of type $(\ell, k)$ can be regarded as a multilinear function that takes $k$ vector fields and returns a linear combiation of tensor products of $\ell$ vector fields. 

\begin{example}
    Let $ F : \mathfrak{X} (M) \times \ldots \times \mathfrak{X} (M)  \to C^{\infty} (M) $ be a function that takes $k$ vector fields and returns a function.  It is a $C^{\infty}(M)$-module morphism if it is multilinear and for any $h \in C^{\infty } (M) $ and any vector fields $X_1, \ldots , X_k \in \mathfrak{X}(M)$, one has
    \begin{align*}
     F  (h X _1,   \ldots , X_k) & = F(X_1, hX_2, \ldots , X_k) \\
     & = \ldots  \\
     & = F(X_1, \ldots , hX_k ) \\
      & = h F (X_1, \ldots , X_k) .   
    \end{align*}
    By the above discussion, this happens if and only if $F$ is associated to a tensor field $\tilde{F} \in \mathcal{T}^k_0(M)$. 
\end{example}

A more direct way of seeing this is the following: for finite-dimensional vector spaces, we have a natural isomorphism
\[     V^{\ast} \otimes W \cong Hom(V , W) .                   \]
Iterating this, 
\begin{align*}
    V ^{\ast} \otimes \cdots \otimes V^{\ast} & \cong Hom ( V , V^{\ast} \otimes \cdots \otimes  V^{\ast} ) \\
     & \cong Hom ( V ,  Hom ( V , V ^{\ast}  \otimes \cdots \otimes  V^{\ast} ) ) \\
     & \cong Hom ( V ,  Hom ( V , Hom (V ,  V ^{\ast}  \otimes \cdots \otimes  V^{\ast} ) ) ) \\
     & \ldots \\
     & \cong  Hom (V , Hom (V, \ldots , Hom (V, \mathbb{R}) \cdots )) \\
     & \cong Mult_{k} (V , \ldots , V ; \mathbb{R} ) \\
     & \cong ( V \otimes \cdots \otimes V )^{\ast},
\end{align*}
where the last isomorphism comes from the fact that linear maps 
\[ V \otimes \cdots \otimes V  \to \mathbb{R}  \] 
are in one-to-one correspondence with $k$-multilinear maps 
\[  V \times \cdots \times V \to \mathbb{R}. \] 
This show that there is a natural isomorphism 
\[    V ^{\ast} \otimes \cdots \otimes V^{\ast} \cong     ( V \otimes \cdots \otimes V )^{\ast}   .    \]
These isomorphisms pass to  vector bundles, so for any vector bundle $E \to M $ we have
\begin{align*}
    Hom _{C^{\infty}(M)} ( & \mathfrak{X}(M )   \otimes \cdots \otimes   \mathfrak{X}(M) , \Gamma (E) ) \\
    & \cong  \Gamma ( Hom ( TM \otimes \cdots \otimes TM ,  E )) \\
    & \cong \Gamma( ( TM \otimes \cdots \otimes TM ) ^{\ast }  \otimes E ) \\
    & \cong \Gamma ( T^{\ast}M \otimes \cdots \otimes T^{\ast} M \otimes E  ). 
\end{align*}




\end{document}